% === §5.3 Enfoques teóricos que sustentan la propuesta ===
% Redactado por el Investigador — Fase 4
% Cada enfoque se vincula con causa-efecto explícito a un subproblema de §2.1

Los enfoques que se describen a continuación constituyen el arsenal metodológico y
algorítmico con que la propuesta aborda cada uno de los tres subproblemas
identificados. Para cada enfoque se presenta su fundamento matemático-computacional,
se argumenta su pertinencia mediante una relación explícita de causa-efecto con el
subproblema que resuelve, y se demuestra su viabilidad técnica en el contexto del
proyecto.

\medskip

% ===== ENFOQUE 1: Modelado computacional del aprendiz (SP1) =====
\noindent\textbf{Enfoque 1. Modelado computacional multidimensional del aprendiz
mediante NLP y aprendizaje automático supervisado.}
\emph{Fundamento.}
Este enfoque combina técnicas de procesamiento de lenguaje natural (NLP) basadas en
transformers —incluyendo modelos preentrenados como BERT, RoBERTa y sus variantes
multilingües ajustadas al español— con clasificadores supervisados (XGBoost,
\emph{random forests}, \emph{support vector machines}) para construir un modelo
computacional del aprendiz que infiere, a partir de trazas textuales, de código y de
interacción, el nivel de aprendizaje profundo en sus tres dimensiones: motivación
intrínseca, argumentación técnica y resolución de problemas auténticos. Las
representaciones semánticas obtenidas mediante embeddings contextualizados alimentan
una arquitectura de clasificación multietiqueta entrenada sobre un \emph{dataset}
anotado con rúbricas validadas psicométricamente.

\emph{Relación causa-efecto con SP1.}
El subproblema SP1 establece que la ausencia de diagnósticos sistemáticos y
trazables del aprendizaje profundo bloquea cualquier intento de personalización
pedagógica. La \textbf{causa} de este bloqueo es doble: (a) la inexistencia de
representaciones computacionales que capturen las dimensiones cognitivas profundas
—más allá de indicadores superficiales de actividad en plataforma—, y (b) la
dependencia de instrumentos psicométricos exclusivamente manuales que no escalan al
aula digital. El \textbf{efecto} de aplicar este enfoque es la producción de un
modelo del aprendiz que transforma señales multimodales —textos argumentativos,
soluciones de código, trazas de simulación— en un perfil cuantitativo y trazable de
aprendizaje profundo, operable en tiempo real por los agentes del ecosistema. Al
capturar la semántica de las producciones del estudiante mediante embeddings
contextualizados y clasificarlas en las tres dimensiones del constructo, el enfoque
resuelve directamente ambas causas del SP1.

\emph{Viabilidad técnica.}
El equipo investigador cuenta con experiencia previa en procesamiento de señales
textuales mediante NLP y en analítica del aprendizaje con \emph{deep learning}
\cite{escamilla2022artificial}. Los modelos de lenguaje preentrenados en español
(BETO, RoBERTa-BNE) están disponibles públicamente y son desplegables en la
infraestructura Apple Silicon del laboratorio. La recolección de datos se realizará
en asignaturas reales de los programas de ingeniería de la UNAL, con protocolos de
consentimiento informado aprobados, garantizando un \emph{dataset} ecológicamente
válido de al menos 200 estudiantes.

\medskip

% ===== ENFOQUE 2: Arquitectura multiagente con LLMs y MCP (SP2) =====
\noindent\textbf{Enfoque 2. Arquitectura multiagente pedagógicamente fundada con
orquestación basada en LLMs y el Model Context Protocol (MCP).}
\emph{Fundamento.}
La arquitectura propuesta sigue el patrón de sistemas multiagente heterogéneos:
agentes especializados —tutor, evaluador, generador de problemas, simulador de
sistemas eléctricos/electrónicos y asistente de programación— coordinan sus acciones
a través de un bus de orquestación que implementa el Model Context Protocol (MCP).
Cada agente encapsula un LLM ajustado mediante ingeniería de \emph{prompts} y,
cuando se requiere precisión factual en dominios ingenieriles, mecanismos de
generación aumentada por recuperación (RAG) sobre una base de conocimiento
documental curada (libros de texto, \emph{datasheets}, manuales de referencia). El
protocolo MCP estandariza el intercambio de contexto entre agentes y herramientas
externas —intérpretes de Python, simuladores SPICE, entornos de compilación—,
garantizando la seguridad, la trazabilidad y la auditabilidad de cada interacción.

\emph{Relación causa-efecto con SP2.}
El SP2 identifica como limitación central la fragmentación de herramientas de IA
educativa que operan como islas sin un modelo compartido del aprendiz, sin memoria
pedagógica y sin orquestación didáctica. La \textbf{causa} raíz es arquitectónica:
los sistemas actuales son monolíticos (un único LLM conversacional) o inconexos
(múltiples herramientas sin protocolo de interoperabilidad), lo que impide la
trayectoria de aprendizaje asistida por agentes. El \textbf{efecto} de este enfoque
es la construcción de un ecosistema que (a) descompone la tarea educativa compleja
en roles agénticos con responsabilidades pedagógicas bien definidas; (b) comparte un
modelo del aprendiz común que dota de continuidad a la trayectoria de aprendizaje;
(c) gradúa la asistencia mediante \emph{scaffolding} adaptativo; y (d) asegura la
seguridad e interoperabilidad mediante MCP, permitiendo que el ecosistema se
comunique con simuladores, compiladores y entornos de ejecución reales. Ninguna de
estas capacidades está presente en los sistemas actuales, y su integración es la
novedad central de la propuesta.

\emph{Viabilidad técnica.}
El protocolo MCP, desarrollado por Anthropic como estándar abierto, cuenta con SDKs
en Python y TypeScript, documentación extensa y una comunidad activa de
desarrolladores \cite{anthropic2024mcp}. Los LLMs necesarios —tanto propietarios
(GPT-4o, Claude 3.5 Sonnet) como abiertos (Llama 3, Mistral)— están disponibles y
han sido probados por el equipo en proyectos previos. La orquestación se implementará
sobre \emph{frameworks} como LangGraph, que han demostrado su eficacia en la
coordinación de flujos multiagente con estado.

\medskip

% ===== ENFOQUE 3: Scaffolding adaptativo computacional (SP2) =====
\noindent\textbf{Enfoque 3. Scaffolding adaptativo computacional basado en la Zona
de Desarrollo Próximo.}
\emph{Fundamento.}
Este enfoque operacionaliza la ZDP de Vygotsky como una función de decisión
computacional $s: \mathbb{R}^d \to \{1, 2, 3\}$ que, dado el vector de
características del aprendiz $\mathbf{x} \in \mathbb{R}^d$ producido por el modelo
del aprendiz (Enfoque 1), asigna un \emph{nivel de scaffolding} $s(\mathbf{x})$ que
controla tres parámetros de la interacción agéntica: (i) $h \in [0,1]$ —profundidad
heurística de las explicaciones, desde sugerencias mínimas ($h \approx 0$) hasta
tutoriales completos ($h \approx 1$)—; (ii) $d \in [0,1]$ —grado de direccionalidad
en la retroalimentación, desde preguntas socráticas abiertas ($d \approx 0$) hasta
corrección explícita ($d \approx 1$)—; y (iii) $c \in \{1, \dots, 5\}$ —complejidad
de los problemas auténticos presentados, calibrada según la taxonomía SOLO. La
función $s$ se implementa mediante un árbol de decisión calibrado con umbrales
derivados de la validación psicométrica del instrumento de diagnóstico y refinado
iterativamente durante la fase piloto.

\emph{Relación causa-efecto con SP2.}
El SP2 señala que los sistemas de IA educativa actuales no personalizan la
asistencia en función del nivel real de competencia del aprendiz, ofreciendo el mismo
tipo de ayuda independientemente de si el estudiante es novato o avanzado. La
\textbf{causa} es la ausencia de un mecanismo computacional que traduzca el perfil
del aprendiz (producido por el Enfoque 1) en decisiones pedagógicas graduadas. El
\textbf{efecto} del scaffolding adaptativo es que cada interacción del ecosistema se
ajusta dinámicamente: un estudiante con baja competencia en argumentación recibe
preguntas estructuradas y modelado explícito de razonamiento; uno con competencia
alta recibe problemas abiertos y retroalimentación socrática. Esta graduación evita
tanto la sobreasistencia —que genera dependencia— como la subasistencia —que produce
frustración y abandono—, dos modos de fallo bien documentados en la literatura sobre
ITS \cite{hooshyar2024systematic}.

\emph{Viabilidad técnica.}
El enfoque no requiere entrenamiento de modelos adicionales: opera como una capa de
decisión ligera (\emph{lookup table} o árbol de decisión) que consume las salidas del
modelo del aprendiz y configura, vía MCP, los parámetros de comportamiento de los
agentes del ecosistema. La calibración de los umbrales se realizará con datos del
piloto, utilizando análisis de curvas ROC y retroalimentación de expertos en
educación en ingeniería.

\medskip

% ===== ENFOQUE 4: Generación de problemas auténticos y evaluación formativa (SP2) =====
\noindent\textbf{Enfoque 4. Generación de problemas auténticos de ingeniería y
evaluación formativa automatizada mediante RAG, \emph{prompt engineering} y NLP.}
\emph{Fundamento.}
Este enfoque integra dos capacidades complementarias del ecosistema agéntico. Por un
lado, la generación de problemas auténticos —tareas de ingeniería abiertas,
contextualizadas y con múltiples caminos de solución— se implementa mediante una
tubería de generación aumentada por recuperación (RAG): una consulta semántica
recupera problemas canónicos y sus soluciones de una base documental de ingeniería
(libros de texto, exámenes de olimpiadas, \emph{datasheets} de componentes), y un
LLM produce una variante inédita condicionada por el nivel de complejidad $c$
asignado por el \emph{scaffolding} adaptativo y por el dominio de especialidad del
estudiante (eléctrica, electrónica, computación). Por otro lado, la evaluación
formativa automatizada emplea NLP para (a) analizar la calidad argumentativa y la
corrección técnica de las respuestas abiertas de los estudiantes mediante similitud
semántica (similitud coseno sobre embeddings de Sentence-BERT) y detección de
entidades ingenieriles (reconocimiento de entidades nombradas, NER, adaptado al
dominio), y (b) generar retroalimentación formativa estructurada según rúbricas que
distinguen entre error conceptual, error procedimental y omisión.

\emph{Relación causa-efecto con SP2.}
El SP2 identifica la carencia de un ecosistema que articule generación de problemas y
evaluación formativa en un mismo flujo pedagógico. La \textbf{causa} es que las
herramientas actuales tratan la generación y la evaluación como funciones
independientes: un LLM puede generar un problema, pero no tiene acceso al modelo del
aprendiz para calibrar su dificultad, ni a un módulo de evaluación que cierre el
ciclo de retroalimentación. El \textbf{efecto} de integrar ambas funciones en el
ecosistema agéntico —con el agente generador y el agente evaluador compartiendo el
modelo del aprendiz y los parámetros de \emph{scaffolding}— es un ciclo completo de
aprendizaje profundo: problema auténtico calibrado $\to$ respuesta del estudiante
$\to$ evaluación semántica $\to$ retroalimentación formativa graduada $\to$ ajuste
del modelo del aprendiz $\to$ siguiente problema. Este ciclo iterativo constituye la
mecánica central mediante la cual el ecosistema fortalece progresivamente las tres
dimensiones del aprendizaje profundo.

\emph{Viabilidad técnica.}
La generación de problemas mediante RAG se apoya en \emph{frameworks} maduros
(LangChain, LlamaIndex) y en modelos de embeddings eficientes (text-embedding-3,
all-MiniLM-L6-v2) que son ejecutables en infraestructura local. La evaluación
formativa automatizada se beneficia de la extensa investigación en \emph{automated
short answer scoring} (ASAS) y \emph{automated essay scoring} (AES), cuyos métodos
han alcanzado coeficientes de acuerdo con evaluadores humanos —medidos mediante
\emph{quadratic weighted kappa}— superiores a $0.75$ en varios \emph{benchmarks}
\cite{chan2024students}. El ajuste al dominio ingenieril se realizará mediante
\emph{fine-tuning} con un corpus anotado de respuestas de estudiantes de ingeniería
de la UNAL.

\medskip

% ===== ENFOQUE 5: Diseño cuasiexperimental pre-post con grupo control (SP3) =====
\noindent\textbf{Enfoque 5. Diseño cuasiexperimental pre-post con grupo de control
no equivalente y analítica del aprendizaje para la validación del impacto causal.}
\emph{Fundamento.}
Este enfoque adopta el diseño cuasiexperimental como estrategia de validación
empírica del ecosistema agéntico, dado que la asignación aleatoria de estudiantes a
condiciones experimentales no es factible en entornos reales de aula universitaria
sin alterar la dinámica académica. El diseño se estructura en tres momentos: (T0)
medición basal (\emph{pretest}) de las tres dimensiones del aprendizaje profundo
mediante el instrumento psicométrico desarrollado en el OE1; (T1) intervención
diferenciada —el grupo experimental interactúa con el ecosistema agéntico durante un
semestre académico, mientras el grupo de control recibe instrucción tradicional con
acceso a herramientas de IA convencionales—; (T2) medición final (\emph{posttest})
con el mismo instrumento. El análisis estadístico emplea ANOVA de medidas repetidas
con factor inter-sujetos (grupo) e intra-sujetos (tiempo), complementado con el
cálculo del tamaño del efecto ($d$ de Cohen, $\eta^2$ parcial) y con análisis de
covarianza (ANCOVA) para controlar diferencias basales. La analítica del aprendizaje
—\emph{learning analytics}— proporciona una capa complementaria de evidencia
cuantitativa: trazas de interacción con el ecosistema (frecuencia de uso, tiempo en
tarea, tasa de persistencia, patrones de solicitud de ayuda) que permiten modelar la
relación dosis-respuesta entre exposición al ecosistema y ganancia en aprendizaje
profundo.

\emph{Relación causa-efecto con SP3.}
El SP3 diagnostica la carencia de evidencia causal —no meramente correlacional— sobre
el impacto de los ecosistemas agénticos en el aprendizaje profundo. La \textbf{causa}
es metodológica: la mayoría de los estudios en IA educativa emplean diseños
pre-experimentales (un solo grupo, solo \emph{posttest}) o estudios de caso sin grupo
de control, que no permiten descartar explicaciones alternativas (maduración, efecto
Hawthorne, regresión a la media). El \textbf{efecto} del diseño cuasiexperimental
pre-post con grupo control es la capacidad de aislar el efecto atribuible al
ecosistema agéntico, controlando estadísticamente las diferencias basales y los
factores de confusión, y produciendo estimaciones del tamaño del efecto comparables
con la literatura internacional sobre intervenciones educativas efectivas
\cite{freeman2014active}. La combinación con analítica del aprendizaje añade una capa
de validez interna al permitir correlacionar la intensidad de uso del ecosistema con
la magnitud de la ganancia observada.

\emph{Viabilidad técnica.}
La UNAL Sede Manizales ofrece el entorno ideal para este diseño: múltiples grupos
paralelos en asignaturas comunes de ingeniería (Circuitos Eléctricos, Programación
Orientada a Objetos, Sistemas Embebidos) permiten constituir grupos experimental y
control con tamaño muestral suficiente ($n \geq 30$ por grupo) y controlar el efecto
del instructor. El equipo investigador cuenta con experiencia en evaluación de
intervenciones educativas \cite{escamilla2022artificial}, y los protocolos éticos
(incluyendo la equiparación de condiciones al finalizar el estudio para el grupo
control) han sido diseñados conforme a los lineamientos del Comité de Ética de la
Facultad de Ingeniería y Arquitectura de la UNAL.

\medskip

% ===== ENFOQUE 6: Marco TRL adaptado a edtech (SP3) =====
\noindent\textbf{Enfoque 6. Marco de madurez tecnológica TRL adaptado a tecnologías
educativas para la transferencia al sector productivo.}
\emph{Fundamento.}
La escala de \emph{Technology Readiness Level} (TRL), originada en la NASA y
adoptada por la Comisión Europea en su programa Horizonte 2020, clasifica la madurez
de una tecnología en nueve niveles, desde la investigación básica (TRL~1) hasta el
despliegue operativo probado (TRL~9). La presente propuesta adapta esta escala al
dominio de las tecnologías educativas (\emph{edtech}) definiendo criterios
operacionales específicos para cada nivel: TRL~4 —validación de los agentes
individuales en entorno de laboratorio con \emph{datasets} históricos—; TRL~5
—validación del ecosistema integrado en entorno simulado con estudiantes voluntarios—;
TRL~6 —demostración del ecosistema en entorno relevante (aula universitaria real) con
evidencia de impacto en aprendizaje profundo—; y TRL~7 —demostración operativa del
ecosistema en un entorno del sector productivo (empresa de software, centro de
formación técnica) con documentación completa de interoperabilidad y operación. La
verificación de cada nivel se realiza mediante una lista de cotejo estandarizada que
evalúa la completitud funcional, la robustez, la interoperabilidad y la evidencia de
impacto.

\emph{Relación causa-efecto con SP3.}
El SP3 identifica la ausencia de una ruta de transferencia tecnológica que garantice
la madurez y la escalabilidad del ecosistema más allá del laboratorio. La
\textbf{causa} es que la mayoría de los prototipos de IA educativa en ingeniería se
desarrollan y evalúan exclusivamente en condiciones controladas (TRL~3--4), sin un
plan explícito de maduración tecnológica ni criterios verificables de avance. El
\textbf{efecto} del marco TRL adaptado es doble: (a) proporciona una hoja de ruta
verificable y estandarizada que guía el desarrollo del ecosistema desde la validación
en laboratorio (punto de partida del equipo) hasta la demostración operativa en el
sector productivo, y (b) genera evidencia documentada —en forma de listas de cotejo
cumplimentadas, informes de desempeño en entorno relevante y testimonios de actores
externos— que reduce la incertidumbre para los potenciales adoptantes (universidades,
centros de formación, empresas) y facilita la transferencia efectiva de la
tecnología.

\emph{Viabilidad técnica.}
El punto de partida del equipo investigador en modelado con \emph{machine learning},
NLP y despliegue de sistemas de IA \cite{escamilla2022artificial} sitúa la línea base
en TRL~3--4 para los componentes individuales. La financiación y la duración de la
convocatoria (18 meses de ejecución más 2 de productos) permiten avanzar dos niveles
en la escala, alcanzando TRL~6 con un plan documentado de proyección a TRL~7. La
alianza con al menos un actor del sector productivo del eje cafetero, prevista en el
OE3, proporciona el entorno operativo necesario para la demostración TRL~7.
